\subsection{CONVOLUTION OF POINT DISTRIBUTIONS ON A LIE GROUP}

DEFINITION 1. Let G be a Lie group, g and \(g'\) two points of G and let \(t \in \mathrm{T}_{g}^{(\infty)}(\mathrm{G})\), \(t' \in \mathrm{T}_{g'}^{(\infty)}(\mathrm{G})\) be two point distributions at g and \(g'\) on G (Differentiable and Analytic Manifolds, R, 13.2.1). The convolution product of t and \(t'\) , denoted by \(t * t'\) , is the image of \(t \otimes t'\) under the mapping \((h, h') \mapsto hh'\) of G × G into G (Differentiable and Analytic Manifolds, R, 13.2.3).

PROPOSITION 1. (i) If \(t \in \mathrm{T}_g^{(s)}(\mathrm{G})\) and \(t' \in \mathrm{T}_{g'}^{(s')}(\mathrm{G})\) , then \(t * t' \in \mathrm{T}_{gg'}^{(s + s')}(\mathrm{G})\) .

\begin{enumerate}
\def\labelenumi{(\roman{enumi})}
\setcounter{enumi}{1}
\tightlist
\item
  If \(t\) or \(t'\) has no constant term, \(t * t'\) has no constant term.
\end{enumerate}

\[
\varepsilon_ {g} * \varepsilon_ {g ^ {\prime}} = \varepsilon_ {g g ^ {\prime}}.
\]

\begin{enumerate}
\def\labelenumi{(\roman{enumi})}
\setcounter{enumi}{3}
\tightlist
\item
  Let \(t \in \mathrm{T}_g^{(s)}(\mathrm{G})\) , \(t' \in \mathrm{T}_{g'}^{(s')}(\mathrm{G})\) and let \(f\) be a function of class \(\mathbf{C}^{s + s'}\) in an open neighbourhood of \(gg'\) with values in a Hausdorff polynormed space. Then
\end{enumerate}

\[
\begin{array}{c} \langle t * t ^ {\prime}, f \rangle = \langle t ^ {\prime}, h ^ {\prime} \mapsto \langle t, h \mapsto f (h h ^ {\prime}) \rangle \rangle \\ = \langle t, h \mapsto \langle t ^ {\prime}, h ^ {\prime} \mapsto f (h h ^ {\prime}) \rangle \rangle . \end{array}
\]

This follows from Differentiable and Analytic Manifolds, R, 13.4.1, 13.2.3 and 13.4.4.

Suppose that K = R or C and that G is finite-dimensional. Then G is locally compact. If t, \(t'\) are point measures, the definition of \(t \star t'\) agrees with that of Integration, Chapter VIII, § 1. We shall see later that the convolution product of measures and that of point distributions are two special cases of the convolution product of distributions which are not necessarily point distributions.

Let \(\mathcal{T}^{(\infty)}(\mathrm{G})\) be the direct sum of the \(\mathrm{T}_g^{(\infty)}(\mathrm{G})\) for \(g\in \mathrm{G}\) (cf.~Differentiable and Analytic Manifolds, R, 13.6.1). We define the convolution product of \(\mathcal{T}^{(\infty)}(\mathrm{G})\) as the bilinear mapping of \(\mathcal{T}^{(\infty)}(\mathrm{G})\times \mathcal{T}^{(\infty)}(\mathrm{G})\) into \(\mathcal{T}^{(\infty)}(\mathrm{G})\) extending the convolution product of Definition 1. We also denote it by \(*\) . Thus \(\mathcal{T}^{(\infty)}(\mathrm{G})\) has an algebra structure filtered by the \(\mathcal{T}^{(s)}(\mathrm{G})\) . The subalgebra \(\mathcal{T}^{(0)}(\mathrm{G}) = \bigoplus_{g\in \Omega}\mathrm{T}_g^{(0)}(\mathrm{G})\) is identified with the algebra \(\mathrm{K}^{\mathrm{(G)}}\) of the group G over K.

PROPOSITION 2. The algebra \(\mathcal{T}^{(\infty)}(\mathrm{G})\) is associative. It is commutative if and only if \(\mathrm{G}\) is commutative.

Let \(t \in \mathcal{T}^{(\infty)}(G)\) , \(t' \in \mathcal{T}^{(\infty)}(G)\) , \(t'' \in \mathcal{T}^{(\infty)}(G)\) . Then \(t * (t' * t'')\) is the image of \(t \otimes t' \otimes t''\) under the mapping \((g, g', g'') \mapsto g(g'g'')\) of \(G \times G \times G\) into G and \((t * t') * t''\) is the image of \(t \otimes t' \otimes t''\) under the mapping \((g, g', g'') \mapsto (gg')g''\) of \(G \times G \times G\) into G. Hence \((t * t') * t'' = t * (t' * t'')\) . It is seen similarly that, if G is commutative, \(t * t' = t' * t\) . If the convolution product is commutative, G is commutative by Proposition 1 (iii).

PROPOSITION 3. If \(t \in \mathcal{T}^{(\infty)}(G)\) and \(g \in G\) , then \(\gamma(g)_*t = \varepsilon_g * t\) , \(\delta(g)_*t = t * \varepsilon_{g^{-1}}\) , (Int \(g)_*t = \varepsilon_g * t * \varepsilon_{g^{-1}}\) . In particular, \(\varepsilon_e\) is the unit element of \(\mathcal{T}^{(\infty)}(G)\) .

Consider the diagram

\[
\mathrm{G} \xrightarrow {\Phi} \mathrm{G} \times \mathrm{G} \xrightarrow {\Psi} \mathrm{G}
\]

where \(\phi\) is the mapping \(h\mapsto (g,h)\) and \(\psi\) is the mapping \((h^{\prime},h)\mapsto h^{\prime}h\) . Then \(\gamma (g) = \psi \circ \phi\) and hence \(\gamma (g)_*t = \psi_*(\phi_*(t))\) . But \(\phi_{*}(t) = \varepsilon_{g}\otimes t\) and hence \(\psi_{*}(\phi_{*}(t)) = \varepsilon_{g}*t\) . The argument is similar for \(\delta (g)_*t\) . Finally,

\[
\operatorname{Int} g = \gamma (g) \circ \delta (g)
\]

and hence \((\operatorname{Int} g)_* = \gamma(g)_* \circ \delta(g)_*\) .

It is therefore seen that, for \(t \in \mathrm{T}(\mathrm{G})\) , \(\varepsilon_g * t\) and \(t * \varepsilon_g\) are equal to \(gt\) and \(tg\) calculated in the group \(\mathrm{T}(\mathrm{G})\) (§ 2, no. 2). But it should be noted that, for \(t\) , \(t'\) in \(\mathrm{T}(\mathrm{G})\) , the product \(tt'\) in the sense of § 2 is in general different from \(t * t'\) .

DEFINITION 2. Let \(\mathbf{G}\) be a Lie group. The subalgebra of \(\mathcal{T}^{(\infty)}(\mathbf{G})\) consisting of the distributions with support contained in \(e\) is denoted by \(\mathbf{U}(\mathbf{G})\) .

This algebra is filtered by the subspaces

\[
\mathrm{U} _ {s} (\mathrm{G}) = \mathrm{U} (\mathrm{G}) \cap \mathcal {T} ^ {(s)} (\mathrm{G}) = \mathrm{T} _ {e} ^ {(s)} (\mathrm{G}).
\]

We write \(\mathrm{U}^{+}(\mathrm{G}) = \mathrm{T}_{e}^{(\infty)+}(\mathrm{G}),\mathrm{U}_{s}^{+}(\mathrm{G}) = \mathrm{U}^{+}(\mathrm{G})\cap \mathrm{U}_{s}(\mathrm{G})\) (cf.~Differentiable and Analytic Manifolds, R, 13.2.1). Recall that \(\mathrm{U}_0(\mathrm{G})\) is identified with K and \(\mathrm{U}_1^+\) (G) with the tangent space \(\mathrm{T_e(G)}\) . In \(\mathrm{U}(G)\) , \(\mathrm{U}^{+}(\mathrm{G})\) is a two-sided ideal supplementary to \(\mathrm{U}_0(\mathrm{G})\) .

Example. Let E be a complete normable space considered as a Lie group. Then the vector space U(E) is canonically identified with the vector space TS(E) (Differentiable and Analytic Manifolds, R, 13.2.4). Let \(m: \mathrm{E} \times \mathrm{E} \to \mathrm{E}\) be addition on E. Then

\[
m _ {*} \colon \mathsf {T S} (\mathrm{E} \times \mathrm{E}) \rightarrow \mathsf {T S} (\mathrm{E})
\]

is equal to \(\mathsf{TS}(m)\) (Differentiable and Analytic Manifolds, R, 13.2.4). For \(t\) , \(t'\) in \(\mathrm{U(E)} = \mathrm{TS(E)}\) , the image \(t * t'\) of the symmetric tensor product \(t \otimes t'\) under \(m_*\) is therefore \(\mathsf{TS}(m)(t \otimes t')\) . By Algebra, Chapter IV, § 5, no. 6, Proposition 7, this image is just the product \(tt'\) in the algebra \(\mathsf{TS(E)}\) . Thus the algebra \(\mathrm{U(E)}\) is identified with the algebra \(\mathsf{TS(E)}\) .

PROPOSITION 4. Consider the bilinear mapping \((u,v)\mapsto u*v\) (resp. \((u,v)\mapsto v*u\) ) of \(\mathrm{U}(\mathrm{G})\otimes \mathbf{K}^{(\mathrm{G})}\) into \(\mathcal{T}^{(\infty)}(\mathrm{G})\) . The corresponding linear mapping of \(\mathrm{U}(\mathrm{G})\otimes \mathbf{K}^{(\mathrm{G})}\) into \(\mathcal{T}^{(\infty)}(\mathrm{G})\) is a vector space isomorphism.

\(K^{(G)}\) is the direct sum of the \(K\varepsilon_{x}\) for \(x\in G\) . On the other hand, the mapping \(u\mapsto u*\varepsilon_{g}\) (resp. \(u\mapsto\varepsilon_{g}*u\) ) is an isomorphism of the vector space \(U(G)=\mathcal{T}_{e}^{(\infty)}(G)\) onto the vector space \(\mathcal{T}_{g}^{(\infty)}(G)\) by Proposition 3. Finally, \(\mathcal{T}^{(\infty)}(G)\) is the direct sum of the \(T_{g}^{(\infty)}(G)\) for \(g\in G\) .

Let X be a manifold of class \(C^{r}\) ( \(r \geqslant \infty\) ) and \(x \in X\) . We have defined (Differentiable and Analytic Manifolds, R, 13.3.1) a canonical filtration on the vector space \(\mathcal{T}_{x}^{(\infty)}(\mathbf{X})\) and a canonical isomorphism \(i_{X,x}\) of the associated graded vector space onto the graded vector space \(\mathsf{TS}(\mathrm{T}_{x}(\mathbf{X}))\) . In particular, let \(\mathrm{T}_{e}(\mathrm{G}) = L\) ; then \(i_{G,e}\) is an isomorphism of the graded vector space gr \(\mathrm{U}(\mathrm{G})\) onto the graded vector space TS(L). But \(\mathrm{U}(\mathrm{G})\) is a filtered algebra, from which we obtain a graded algebra structure on gr \(\mathrm{U}(\mathrm{G})\) .

PROPOSITION 5. The isomorphism \(i_{G,e} \colon \operatorname{gr} U(G) \to TS(L)\) is an algebra isomorphism.

Let \(p\) be the mapping \((t, t') \mapsto t \otimes t'\) of U(G) × U(G) into U(G × G). Let \(c\) be the mapping \((t, t') \mapsto t * t'\) of U(G) × U(G) into U(G). Let m be the mapping \((g, g') \mapsto gg'\) of G × G into G. Then by Definition 1

\begin{enumerate}
\def\labelenumi{(\arabic{enumi})}
\tightlist
\item
\end{enumerate}

\[
c = m _ {*} \circ p.
\]

Consider the diagram

\[
\begin{array}{c} \operatorname{gr} \mathrm{U(G)} \times \operatorname{gr} \mathrm{U(G)} \xrightarrow {\operatorname{gr} (p)} \operatorname{gr} \mathrm{U(G \times G)} \xrightarrow {\operatorname{gr} (m _ {*})} \operatorname{gr} \mathrm{U(G)} \\ i _ {\mathrm{G}, e} \times i _ {\mathrm{G}, e} \Biggl \downarrow \\ \mathsf {T S} (\mathrm{L}) \times \mathsf {T S} (\mathrm{L}) \xrightarrow {q} \mathsf {T S} (\mathrm{L} \times \mathrm{L}) \xrightarrow {\mathsf {T S} (\mathrm{T} (m))} \mathsf {T S} (\mathrm{L}) \end{array}
\]

where q is the mapping derived from the canonical isomorphism of TS(L) × TS(L) onto TS(L × L). By Differentiable and Analytic Manifolds, R, 13.4.6 and 13.3.5, the two squares of the diagram are commutative. Hence by (1) the diagram

\[
\begin{array}{c} \operatorname{gr} U (G) \times \operatorname{gr} U (G) \xrightarrow {\operatorname{gr} (c)} \operatorname{gr} U (G) \\ ^ {i _ {G, e} \times i _ {G, e}} \Biggl \downarrow \qquad \qquad \qquad \qquad \qquad \qquad \qquad \Biggl \downarrow^ {i _ {G, e}} \\ T S (L) \times T S (L) \xrightarrow {T S (T (m)) \circ q} T S (L) \end{array}
\]

is commutative. Now \(\mathbf{T}(m): \mathbf{L} \times \mathbf{L} \to \mathbf{L}\) maps \((x, y)\) to \(x + y\) (§ 2, no. 1, Proposition 2 (ii)). By Algebra, Chapter IV, § 5, no. 6, Proposition 7, \(\mathbf{TS}(\mathbf{T}(m)) \circ q\) is therefore the multiplication of the algebra \(\mathbf{TS}(\mathbf{L})\) .

